\section{Meeting Details - 03, Feb 18}
\label{chap:Meeting_03}
Today's meeting will take place in Small Hall Room No. 038. The agenda and detailed minutes of the meeting are written below.
\\ \textbf{Attendes:}
\textcolor{magenta}{\textbf{Pravin Sharma} and \textbf{Emmanuel Sampson}}
\\ \textbf{Type of Meeting:}
\textcolor{blue}{\textbf{In-Person} }
\subsection*{Meeting Agenda}
\begin{enumerate}
    \item Ask Emmanuel about his time availability.
    \item Regarding the deadline, we want this project completed.
    \item Work division for the calculation and data interpretation.
    \item Ask him to add a couple of papers, make separate folders if needed for detailed information, and workflow.
    \item Suggestions to look for the system-related papers, and add them to the library.
    \item Meet regularly on a weekly basis, Zoom/in-person, whichever is feasible.
    \item Ask him to scan about FeRh as well as another candidate for exchange bias related work. \xmark
\end{enumerate}

\subsection*{Meeting Minutes}
\begin{enumerate}
    \item Emmanuel said he can work 2 hours/week on the project.
    \item I clearly told him, if you feel like it's your project, then it won't make the change, but if you really think we can pull this together, then come on and be responsible for the task you are given within the provided deadline.
    \item I told him specifically that he should be focused only on the calculation of hysteresis and do the literature survey related to the parameters of $\mathrm{CoFe/FeO}$ and $\mathrm{CoFe/CoO}$, as these are the materials that have high chances of formation at the interface of $\mathrm{CoFe/MgO}$ due to the oxygen diffusing at the interface.
    \item We agreed to complete the project and have the first draft by mid-May/June, and then go from there. 
    \item I managed to help him understand the parameters and some basics about the processes, which we are looking forward to. 
    \item I have to email him some details about hysteresis and how the study of Exchange-bias (EB) and the Coercive field $(H_c)$ will help us in this work. 
    \item I will share a folder containing the input and .mat files, so he can become familiar with them and work on them clearly.
\end{enumerate}
\stardivider